\documentclass[../../main.tex]{subfiles}
\begin{document}
\section{线性映射}\label{section:appB-linear-maps}

设 $V$ 与 $W$ 是实向量空间。映射 $T:V\to W$ 如果对所有 $v,w\in V$ 与所有标量 $a,b$ 都满足 $T(av+bw)=aTv+bTw$，就称为\emph{线性的}。（由于下文所述线性映射与矩阵乘法之间的密切联系，我们通常把线性映射 $T$ 对向量 $v$ 的作用写作 $Tv$ 而不用括号，除非为分组需要括号。）在特殊情形 $W=\mathbb R$，从 $V$ 到 $\mathbb R$ 的线性映射通常称为 $V$ 上的\emph{线性泛函}。

如果 $T:V\to W$ 是线性映射，$T$ 的\emph{核}或\emph{零空间}记作 $\ker T$ 或 $T^{-1}(0)$，是集合 $\{v\in V:Tv=0\}$；$T$ 的\emph{像}记作 $\operatorname{Im}T$ 或 $T(V)$，是集合 $\{w\in W:w=Tv\text{ 对某个 }v\in V\}$。

一个简单而重要的线性映射例子如下产生。给定子空间 $S\subset V$ 和互补子空间 $T$，存在唯一线性映射 $\pi:V\to S$，定义为
\begin{align*}
\pi(v+w)=v\qquad\text{对 }v\in S,\ w\in T.
\end{align*}
这个映射称为以 $T$ 为核到 $S$ 的\emph{投影}。

如果 $V$ 与 $W$ 是向量空间，双射线性映射 $T:V\to W$ 称为\emph{同构}。在此情形中，存在唯一逆映射 $T^{-1}:W\to V$，而下列计算表明 $T^{-1}$ 也线性：
\begin{align*}
aT^{-1}v+bT^{-1}w=T^{-1}T\big(aT^{-1}v+bT^{-1}w\big)
=T^{-1}\big(aTT^{-1}v+bTT^{-1}w\big)
=T^{-1}(av+bw),
\end{align*}
其中倒数第二步用了 $T$ 的线性性。因此，双射线性映射也称为\emph{可逆的}。如果存在同构 $T:V\to W$，就说 $V$ 与 $W$ \emph{同构}。

\begin{example}\label{example:appB-basis-isomorphism}
设 $V$ 是 $n$ 维实向量空间，$(E_1,\dots,E_n)$ 是 $V$ 的有序基。定义映射 $E:\mathbb R^n\to V$ 为
\begin{align*}
E(x^1,\dots,x^n)=x^1E_1+\cdots+x^nE_n.
\end{align*}
则 $E$ 线性且双射，故是同构，称为由这个基确定的\emph{基同构}。因此，每个 $n$ 维实向量空间都同构于 $\mathbb R^n$。
\end{example}

\begin{exercise}\label{exercise:appB-B-13}
设 $V$ 与 $W$ 是向量空间，$(E_1,\dots,E_n)$ 是 $V$ 的基。对任意 $n$ 个元素 $w_1,\dots,w_n\in W$，证明存在唯一线性映射 $T:V\to W$ 满足 $T(E_i)=w_i$（$i=1,\dots,n$）。
\end{exercise}

\begin{exercise}\label{exercise:appB-B-14}
设 $S:V\to W$ 与 $T:W\to X$ 是线性映射。
\begin{enumerate}[(a)]
\item 证明 $\ker S$ 与 $\operatorname{Im}S$ 分别是 $V$ 与 $W$ 的子空间。
\item 证明 $S$ 单射，当且仅当 $\ker S=\{0\}$。
\item 证明如果 $S$ 是同构，则 $\dim V=\dim W$（其意义为这两个维数要么都无限、要么都有限且相等）。
\item 证明如果 $S$ 与 $T$ 都单射或都满射，则 $T\circ S$ 有相同性质。
\item 证明如果 $T\circ S$ 满射，则 $T$ 满射；举一个例子表明 $S$ 未必满射。
\item 证明如果 $T\circ S$ 单射，则 $S$ 单射；举一个例子表明 $T$ 未必单射。
\end{enumerate}
\end{exercise}

\begin{exercise}\label{exercise:appB-B-15}
假设 $V$ 是向量空间，$S$ 是 $V$ 的子空间，$\pi:V\to V/S$ 表示由 $\pi(v)=v+S$ 定义的投影。
\begin{enumerate}[(a)]
\item 证明 $\pi$ 是满线性映射，核等于 $S$。
\item 给定线性映射 $T:V\to W$，证明存在线性映射 $\widetilde T:V/S\to W$ 使 $\widetilde T\circ\pi=T$，当且仅当 $S\subset\ker T$。
\end{enumerate}
\end{exercise}

如果 $V$ 与 $W$ 是向量空间，映射 $F:V\to W$ 如果可写成形式 $F(v)=w+Tv$（对某个线性映射 $T:V\to W$ 和某个固定 $w\in W$），就称为\emph{仿射映射}。

\begin{exercise}\label{exercise:appB-B-16}
假设 $F:V\to W$ 是仿射映射。证明 $F(V)$ 是 $W$ 的仿射子空间，而集合 $F^{-1}(z)$（$z\in W$）是 $V$ 的平行仿射子空间。
\end{exercise}

\begin{exercise}\label{exercise:appB-B-17}
假设 $V$ 是有限维向量空间。证明 $V$ 的每个仿射子空间都形如 $F^{-1}(z)$（对某个仿射映射 $F:V\to W$ 和某个 $z\in W$）。
\end{exercise}

现设 $V$ 与 $W$ 是有限维向量空间，分别带有有序基 $(E_1,\dots,E_n)$ 和 $(F_1,\dots,F_m)$。如果 $T:V\to W$ 是线性映射，$T$ 关于这些基的\emph{矩阵}是 $m\times n$ 矩阵 $A=(A^i_j)$，其第 $j$ 列由 $TE_j$ 关于基 $(F_i)$ 的分量组成：
\begin{align*}
TE_j=\sum_{i=1}^n A^i_j F_i.
\end{align*}

由线性性，$T$ 对任意向量 $v=\sum_j v^j E_j$ 的作用为
\begin{align*}
T\left(\sum_{j=1}^n v^j E_j\right)=\sum_{i=1}^m\sum_{j=1}^n A^i_j v^j F_i.
\end{align*}
如果我们把向量关于基的分量写成列矩阵，则 $w=Tv$ 的矩阵表示由矩阵乘法给出，或更简洁地写作
\begin{align*}
w^i=\sum_{j=1}^n A^i_j v^j.
\end{align*}
在可能范围内，我们用上标表示矩阵的行指标、下标表示列指标，故 $A^i_j$ 表示第 $i$ 行第 $j$ 列的元素。因而矩阵乘积 $AB$ 的第 $i$ 行第 $j$ 列元素为
\begin{align*}
(AB)^i_j=\sum_{k=1}^n A^i_k B^k_j.
\end{align*}

两个线性映射的复合由它们矩阵的乘积表示。只要定义域和陪域用同一基，$n$ 维向量空间上的恒等映射就由 $n\times n$ 单位矩阵表示，记作 $I_n$；它是主对角线（行号等于列号处）上为 $1$、其余为零的矩阵。

所有 $m\times n$ 实矩阵所成集合 $M(m\times n;\mathbb R)$ 容易看出是维数为 $mn$ 的实向量空间。类似地，由于 $\mathbb C$ 是维数为 $2$ 的实向量空间，$m\times n$ 复矩阵所成集合 $M(m\times n;\mathbb C)$ 是维数为 $2mn$ 的实向量空间。当 $m=n$ 时，我们把 $n\times n$ 实、复方阵空间缩写为 $M(n;\mathbb R)$ 和 $M(n;\mathbb C)$。在此情形中，矩阵乘法给这些空间附加代数结构。如果 $V,W,Z$ 是向量空间，映射 $B:V\times W\to Z$ 如果在固定另一变量时分别对每个变量线性，就称为\emph{双线性的}。

\emph{（$\mathbb R$ 上的）代数}是赋予双线性乘积映射 $V\times V\to V$ 的实向量空间 $V$。如果双线性乘积具有交换性或结合性，就称该代数交换或结合。

\begin{exercise}\label{exercise:appB-B-18}
证明矩阵乘法把 $M(n;\mathbb R)$ 和 $M(n;\mathbb C)$ 都变成 $\mathbb R$ 上的结合代数。证明除非 $n=1$，它们非交换。
\end{exercise}

假设 $A$ 是 $n\times n$ 矩阵。如果存在矩阵 $B$ 使 $AB=BA=I_n$，则称 $A$ \emph{可逆}或\emph{非奇异}；否则称\emph{奇异}。

\begin{exercise}\label{exercise:appB-B-19}
假设 $A$ 是 $n\times n$ 矩阵。证明下列命题。
\begin{enumerate}[(a)]
\item 如果 $A$ 非奇异，则存在唯一 $n\times n$ 矩阵 $B$ 使 $AB=BA=I_n$。这个矩阵记作 $A^{-1}$，称为 $A$ 的\emph{逆}。
\item 如果 $A$ 是线性映射 $T:V\to W$ 关于 $V$ 和 $W$ 某基的矩阵，则 $T$ 可逆，当且仅当 $A$ 可逆；此时 $A^{-1}$ 是 $T^{-1}$ 关于同一组基的矩阵。
\item 如果 $B$ 是 $n\times n$ 矩阵使 $AB=I_n$ 或 $BA=I_n$，则 $A$ 非奇异且 $B=A^{-1}$。
\end{enumerate}
\end{exercise}

由于 $\mathbb R^n$ 配备标准基 $(e_i)$，我们可无歧义地把从 $\mathbb R^n$ 到 $\mathbb R^m$ 的线性映射与 $m\times n$ 实矩阵等同起来，且我们常不加说明地这样做。

\subsection*{基变换}

本书中我们常需关心当更换基时各种对象如何变换。假设 $(E_i)$ 与 $(\widetilde E_i)$ 是有限维实向量空间 $V$ 的两个基。每个基可由另一个基唯一表示，故存在可逆矩阵 $B$，称为两基间的\emph{过渡矩阵}，使
\begin{align}\label{align:appB-transition}
E_i=\sum_{j=1}^n B^j_i\widetilde E_j,\qquad
\widetilde E_j=\sum_{i=1}^n (B^{-1})^i_j E_i.
\end{align}

现设 $V$ 与 $W$ 是有限维向量空间，$T:V\to W$ 是线性映射。关于定义域 $V$ 的基 $(E_i)$ 和陪域 $W$ 的基 $(F_j)$，映射 $T$ 由某矩阵 $A=(A^i_j)$ 表示。如果 $(\widetilde E_i)$ 和 $(\widetilde F_j)$ 分别是 $V$ 和 $W$ 基的任意其他选取，令 $B$ 和 $C$ 表示满足(??)和
\begin{align*}
F_i=\sum_{j=1}^m C^j_i\widetilde F_j,\qquad
\widetilde F_j=\sum_{i=1}^m (C^{-1})^i_j F_i
\end{align*}
的过渡矩阵。直接计算表明，关于新基表示 $T$ 的矩阵 $\widetilde A$ 与 $A$ 的关系为
\begin{align*}
\widetilde A^i_j=\sum_{k,l} C^i_l A^l_k (B^{-1})^k_j,
\end{align*}
或用矩阵记号，
\begin{align*}
\widetilde A=CAB^{-1}.
\end{align*}

特别地，如果 $T$ 是从 $V$ 到自身的映射，我们通常对定义域和陪域用同一基。在此情形中，若 $A$ 表示 $T$ 关于 $(E_i)$ 的矩阵，$\widetilde A$ 是它关于 $(\widetilde E_i)$ 的矩阵，则
\begin{align}\label{align:appB-similarity}
\widetilde A=BAB^{-1}.
\end{align}

如果 $V$ 与 $W$ 是实向量空间，从 $V$ 到 $W$ 的线性映射所成集合 $\mathcal L(V;W)$ 在运算
\begin{align*}
(S+T)v=Sv+Tv,\qquad (cT)v=c(Tv)
\end{align*}
下是实向量空间。如果 $\dim V=n$、$\dim W=m$，则 $V$ 和 $W$ 基的每种选取给出一个映射 $\mathcal L(V;W)\to M(m\times n;\mathbb R)$，把每个线性映射映为它关于所选基的矩阵。这个映射容易看出是线性双射，故 $\dim\mathcal L(V;W)=\dim M(m\times n;\mathbb R)=mn$。

如果 $T:V\to W$ 是有限维空间之间的线性映射，$\operatorname{Im}T$ 的维数称为 $T$ 的\emph{秩}，$\ker T$ 的维数称为它的\emph{零化度}。下一个定理表明：在基的选取范围内，一个线性映射完全由其秩连同定义域与陪域的维数决定。

\begin{theorem}[线性映射的典范形式]\label{theorem:appB-linear-map-canonical}
假设 $V$ 与 $W$ 是有限维向量空间，$T:V\to W$ 是秩为 $r$ 的线性映射。则存在 $V$ 和 $W$ 的基，使 $T$ 关于它们有下列（分块形式的）矩阵表示：
\begin{align*}
\begin{pmatrix}I_r&0\\0&0\end{pmatrix}.
\end{align*}
\end{theorem}

\begin{proof}
为 $\operatorname{Im}T$ 选基 $(F_1,\dots,F_r)$，为 $\ker T$ 选基 $(K_1,\dots,K_k)$。把 $(F_j)$ 任意延拓为 $W$ 的基 $(F_1,\dots,F_m)$。由像的定义，存在向量 $E_1,\dots,E_r\in V$ 使 $TE_i=F_i$（$i=1,\dots,r$）。我们将证明 $(E_1,\dots,E_r,K_1,\dots,K_k)$ 是 $V$ 的基；一旦知道此点，容易推出 $T$ 有所需矩阵表示。

先设 $\sum_i a_iE_i+\sum_j b_jK_j=0$。对此方程作用 $T$ 得 $\sum_{i=1}^r a_iF_i=0$，这蕴含所有系数 $a_i$ 为零。然后由于 $K_j$ 线性无关，也推出所有 $b_j$ 为零。因此 $(r+k)$ 元组 $(E_1,\dots,E_r,K_1,\dots,K_k)$ 线性无关。

为证这些向量张成 $V$，任取 $v\in V$。我们可把 $Tv\in\operatorname{Im}T$ 表为 $(F_1,\dots,F_r)$ 的线性组合：
\begin{align*}
Tv=\sum_{i=1}^r c_iF_i.
\end{align*}
若令 $w=\sum_i c_iE_i\in V$，则 $Tw=Tv$，故 $z=v-w\in\ker T$。写 $z=\sum_j d^jK_j$，得
\begin{align*}
v=w+z=\sum_{i=1}^r c_iE_i+\sum_{j=1}^k d^jK_j,
\end{align*}
故 $(r+k)$ 元组 $(E_1,\dots,E_r,K_1,\dots,K_k)$ 确实张成 $V$。
\end{proof}

这个定理说：通过在定义域和陪域中适当选取基，每个线性映射都可化成特别好的对角形式。然而重要的是注意定理\emph{没有}说什么：如果 $T:V\to V$ 是从有限维向量空间到自身的线性映射，则未必能在定义域和陪域同一基下用对角矩阵表示 $T$。

下一个结果在线性代数应用于光滑流形理论时处于核心地位；它是前一定理证明的一个推论。

\begin{corollary}[秩--零维定律]\label{corollary:appB-rank-nullity}
假设 $T:V\to W$ 是有限维向量空间之间的线性映射。则
\begin{align*}
\dim V=\operatorname{rank}T+\operatorname{nullity}T
=\dim(\operatorname{Im}T)+\dim(\ker T).
\end{align*}
\end{corollary}

\begin{proof}
前一证明表明 $V$ 有一个由 $k+r$ 个元素组成的基，其中 $k=\dim(\ker T)$、$r=\dim(\operatorname{Im}T)$。
\end{proof}

\begin{exercise}\label{exercise:appB-B-22}
假设 $V,W,X$ 是有限维向量空间，$S:V\to W$ 与 $T:W\to X$ 是线性映射。证明下列命题。
\begin{enumerate}[(a)]
\item $\operatorname{rank}S\leqslant\dim V$，等号当且仅当 $S$ 单射。
\item $\operatorname{rank}S\leqslant\dim W$，等号当且仅当 $S$ 满射。
\item 如果 $\dim V=\dim W$ 且 $S$ 单射或满射，则它是同构。
\item $\operatorname{rank}(T\circ S)\leqslant\operatorname{rank}S$，等号当且仅当 $\operatorname{Im}S\cap\ker T=\{0\}$。
\item $\operatorname{rank}(T\circ S)\leqslant\operatorname{rank}T$，等号当且仅当 $\operatorname{Im}S+\ker T=W$。
\item 如果 $S$ 是同构，则 $\operatorname{rank}(T\circ S)=\operatorname{rank}T$。
\item 如果 $T$ 是同构，则 $\operatorname{rank}(T\circ S)=\operatorname{rank}S$。
\end{enumerate}
\end{exercise}

设 $A$ 是 $m\times n$ 矩阵。$A$ 的\emph{转置}是交换 $A$ 的行与列所得的 $n\times m$ 矩阵 $A^T$：$(A^T)^j_i=A^i_j$。方阵 $A$ 如果 $A=A^T$ 就称\emph{对称}，如果 $A=-A^T$ 就称\emph{反称}。

\begin{exercise}\label{exercise:appB-B-23}
证明如果 $A$ 与 $B$ 分别是 $m\times n$ 和 $n\times k$ 矩阵，则 $(AB)^T=B^TA^T$。
\end{exercise}

$m\times n$ 矩阵 $A$ 的\emph{秩}定义为相应从 $\mathbb R^n$ 到 $\mathbb R^m$ 的线性映射的秩。由于 $A$ 的列（看作 $\mathbb R^m$ 中向量）是该线性映射下标准基向量的像，$A$ 的秩也可看作其列张成的维数，有时称为其\emph{列秩}。类似地，定义 $A$ 的\emph{行秩}为其行张成的维数（行同样看作 $\mathbb R^n$ 中向量）。

\begin{proposition}\label{proposition:appB-row-column-rank}
矩阵的行秩等于其列秩。
\end{proposition}

\begin{proof}
设 $A$ 是 $m\times n$ 矩阵。由于 $A$ 的行秩等于 $A^T$ 的列秩，我们只需证 $\operatorname{rank}A=\operatorname{rank}A^T$。

设 $A$ 的（列）秩为 $k$。作为从 $\mathbb R^n$ 到 $\mathbb R^m$ 的线性映射，$A$ 如下通过 $\operatorname{Im}A$ 分解：$\mathbb R^n\to\operatorname{Im}A\hookrightarrow\mathbb R^m$，其中 $\widetilde A$ 只是把 $A$ 的陪域限制到 $\operatorname{Im}A$ 的映射，$\iota$ 是 $\operatorname{Im}A$ 到 $\mathbb R^m$ 的包含。为 $k$ 维子空间 $\operatorname{Im}A$ 选基，我们可把它写成矩阵方程 $A=BC$，其中 $B,C$ 分别是 $\iota$ 和 $\widetilde A$ 关于 $\mathbb R^n$、$\mathbb R^m$ 中标准基和 $\operatorname{Im}A$ 中所选基的矩阵。取转置得 $A^T=C^TB^T$，由此 $\operatorname{rank}A^T\leqslant\operatorname{rank}B^T$。由于 $B^T$ 是 $k\times m$ 矩阵，其列秩至多为 $k$，这表明 $\operatorname{rank}A^T\leqslant\operatorname{rank}A$。交换 $A$ 与 $A^T$ 的角色并利用 $(A^T)^T=A$，得出 $\operatorname{rank}A=\operatorname{rank}A^T$。
\end{proof}

假设 $A=(A^i_j)$ 是 $m\times n$ 矩阵。如果我们选取整数 $1\leqslant i_1<\cdots<i_k\leqslant m$ 和 $1\leqslant j_1<\cdots<j_l\leqslant n$，就得到一个 $k\times l$ 矩阵，其第 $p$ 行第 $q$ 列元素为 $A^{i_p}_{j_q}$。这样的矩阵称为 $A$ 的一个\emph{子矩阵}。考察子矩阵给出检查矩阵秩的一个方便判别法。

\begin{proposition}\label{proposition:appB-submatrix-rank}
假设 $A$ 是 $m\times n$ 矩阵。则 $\operatorname{rank}A\geqslant k$，当且仅当 $A$ 有某个 $k\times k$ 子矩阵非奇异。
\end{proposition}

\begin{proof}
由定义，$\operatorname{rank}A\geqslant k$，当且仅当 $A$ 至少有 $k$ 个线性无关列，这等价于 $A$ 有某个秩为 $k$ 的 $m\times k$ 子矩阵。但由命题 B.24，$m\times k$ 子矩阵秩为 $k$，当且仅当它有 $k$ 个线性无关行。故 $A$ 秩至少为 $k$，当且仅当它有一个含 $k$ 个线性无关行的 $m\times k$ 子矩阵，当且仅当它有一个非奇异 $k\times k$ 子矩阵。
\end{proof}
\end{document}
